\subsection*{NT-E17--E18. A quadratic-Dirichlet input for CM recognition}

Write $A_{\rm quad}$ for RH for $\zeta_{\mathbf Q}$ and
for every primitive quadratic Dirichlet $L$-function.
This is a subfamily of the user-designated input $A_D$.
No Hecke function over a number field of degree greater
than one needs to be assumed separately below.
For a nonzero integer polynomial $H$, put
\[
 h=\deg H,\qquad
 \mathcal H=\max\{1,\log\max_i|H_i|\}.
\]
The height is padded by one to include polynomials
whose nonzero coefficients are units.

\paragraph{Exact classical inputs and scope.}
We use Cremona--Sutherland,
\emph{Computing the endomorphism ring of an elliptic
curve over a number field}, Contemporary Mathematics~796
(2024), 75--102, in author version
arXiv:2301.11169v4 (15 December 2023).
All page numbers in this subsection refer to that
author version.

Its Proposition~19, pp.~9--10, proves under GRH that
every imaginary quadratic order discriminant $D<0$
of class number $h$ satisfies
\[
 |D|\le
 \left(\frac{12e^{2\gamma}}{\pi}\right)^2
 h^2\bigl(\log\log(h+1)+4\bigr)^4.
 \tag{E17.1}
\]
The analytic input in this proof is
Languasco--Trudgian, \emph{Uniform effective estimates
for $|L(1,\chi)|$}, Corollary~1, author p.~2,
applied to the primitive quadratic character of the
fundamental discriminant.
The remaining reduction from maximal to arbitrary
orders uses the class-number formula (2.1) and an
unconditional Euler-totient estimate.
Thus $A_{\rm quad}$ supplies the analytic hypothesis
of (E17.1). The cited small-discriminant
tabulations and numerical range checks in that
published proposition are retained as classical
inputs; they were not rerun here.

Let $\Phi_\ell(X,Y)\in\mathbf Z[X,Y]$ be the classical
modular polynomial for cyclic $\ell$-isogenies.
For monic irreducible $H\in\mathbf Z[X]$ and a root $j$,
define
\[
 m_\ell(j)=
 \operatorname{ord}_H
    \operatorname{Res}_Y\bigl(H(Y),\Phi_\ell(X,Y)\bigr).
 \tag{E17.2}
\]
Cremona--Sutherland equation (2.2), p.~8, gives
this equality with the number of conjugates joined
to $j$ by a cyclic $\ell$-isogeny, with multiplicity.
Their Corollary~14, p.~8, gives $m_\ell(j)=2$
for a CM $j$ of order discriminant $D$ whenever
$\ell\nmid D$ splits in $\mathbf Q(\sqrt D)$.
Their Proposition~16, p.~9, gives
$m_\ell(j)\le1$ for every prime $\ell$ if $j$ is not CM.
These are classical algebraic statements, independent
of GRH. Their supplied derivations are read here;
the upstream theory of isogeny volcanoes and
$K$-curves remains cited.

The published Algorithm~1 has conditional correctness
in Theorem~24, p.~14, and deterministic complexity
in Theorem~27 and its full proof, p.~14.
For that particular algorithm the published time
and space bounds are respectively
$(h^2\mathcal H)^{1+o(1)}$ and
$(h\mathcal H)^{1+o(1)}$.
For our modified test we assert the weaker
polynomial bound, with its exact-arithmetic
verification supplied below.
We use the deterministic polynomial-time computation
of $\Phi_\ell$ explicitly cited in the proof
of Theorem~27.
A direct primary locator is
Blake--Csirik--Rubinstein--Seroussi,
\emph{On the Computation of Modular Polynomials
for Elliptic Curves}, \S2, pp.~4--7:
the exact power-series and CRT algorithm is
given there, with the bit-complexity discussion
and coefficient bound (9) on p.~7.
The modified test below examines all the bounded
degrees, so it does not require the source's
least-ordinary-prime selection or its finite list
of abelian discriminants.

\paragraph{NT-E17: a bounded algebraic CM certificate under $A_D$.}
Let $H\in\mathbf Z[X]$ be monic and irreducible,
of degree $h\ge1$, and set
\[
 b_h=\operatorname{bitlength}(h+1),\qquad
                     P_h=(72b_h+264)^2.
 \tag{E17.3}
\]
Under $A_{\rm quad}$, $H$ is a Hilbert class
polynomial for an imaginary quadratic order if
and only if there is a prime $5\le\ell\le P_h$ with
\[
 H(X)^2\ \bigm|\ 
       \operatorname{Res}_Y\bigl(H(Y),\Phi_\ell(X,Y)\bigr)
                           \quad\hbox{in }\mathbf Z[X].
 \tag{E17.4}
\]
Scanning these primes and performing the exact
divisibility tests is deterministic and takes
polynomial time in $h\mathcal H$.
A positive answer, accompanied by $\ell$ and
the exact divisibility identity, is valid
unconditionally. The universal correctness of
a negative answer uses $A_{\rm quad}$.
Irreducibility is part of the stated input promise.

\begin{proof}
If (E17.4) holds, (E17.2) gives $m_\ell(j)\ge2$.
The cited non-CM bound therefore forces $j$ to be CM.
Its monic irreducible polynomial is then exactly
the Hilbert class polynomial of its endomorphism
order. This implication is unconditional.

For the converse suppose $H=H_D$.
First we replace (E17.1) by an elementary bound
whose constants can be used without computing
$\gamma$, $\pi$, or logarithms in the algorithm.
For $h\ge1$,
\[
 \log\log(h+1)+4
 \le\log(h+1)+4
 \le h+4\le3(h+1).
\]
Using $\gamma<1$, $e<3$, and $\pi>3$ in (E17.1),
we obtain
\[
 |D|<1296\cdot81\,h^2(h+1)^4
       <2^{17}(h+1)^6.
 \tag{E17.5}
\]
This bound includes $D=-3,-4$ and nonmaximal orders
because (E17.1) does.

Set $m=\operatorname{rad}(6|D|)$ and
\[
 K=\mathbf Q(\sqrt D),\qquad
 T=K(\sqrt m).
\]
The integer $m$ is positive, even, and squarefree.
The field $\mathbf Q(\sqrt m)$ is real quadratic
and $K$ is imaginary quadratic, so $T/\mathbf Q$
is biquadratic of degree four.
It is unramified outside the primes dividing $m$.
Every prime dividing $m$ ramifies already in
$\mathbf Q(\sqrt m)$, since its discriminant is $4m$.
Consequently every unramified prime of $T$ is
prime to $6D$.

The full Euler factorization of $\zeta_T$ is
$\zeta_{\mathbf Q}$ times the three quadratic
Dirichlet functions of its quadratic subfields.
Thus $A_{\rm quad}$ supplies RH for $\zeta_T$
and $\zeta_{\mathbf Q}$.
Apply NT-C03 to the identity class of
$\operatorname{Gal}(T/\mathbf Q)$.
There is a prime $\ell$ splitting completely in $T$,
with $\ell\nmid6D$ and
\[
                 \ell\le(4\log|\operatorname{disc}T|+15)^2.
\]
The Galois different bound proved in NT-E16,
with absolute degree four, gives
\[
 \log|\operatorname{disc}T|
       \le3\log m+4\log4
       \le3\log(6|D|)+4\log4.
 \tag{E17.6}
\]
It follows that
\[
 \begin{aligned}
 4\log|\operatorname{disc}T|+15
   &\le12\log|D|+12\log6+16\log4+15\\
   &<12\log|D|+60\\
   &<72\log(h+1)+204\log2+60\\
   &<72b_h+264.
 \end{aligned}
 \tag{E17.7}
\]
For the second line one can use
$\log6<9/5$ and $\log4<7/5$, which give the
stronger numerical sum $59$.
Both inequalities follow by summing the exponential
series through degree six:
the resulting positive margins over $6$ and $4$
are respectively $42629/1250000$ and
$592609/11250000$.
The last line uses $h+1<2^{b_h}$ and $\log2<1$.
We have proved $5\le\ell\le P_h$, $\ell\nmid D$,
and splitting in $K$.
The cited CM formula gives $m_\ell(j)=2$,
and hence (E17.4).

Finally $P_h=O(\log^2(h+2))$, so both the number
and the sizes of the tested primes are polynomial
in the input length.
The cited deterministic construction of
$\Phi_\ell$ takes polynomial time in $\ell$,
and its output has polynomial bit size.
Here is an elementary polynomial bound for
the rest of the exact computation.
The Sylvester matrix defining the resultant
has size $N=h+\ell+1$, and its entries are
integer polynomials in $X$ of degree at most
$\ell+1$ with polynomially bounded coefficient
bit lengths.
Its determinant has degree at most
$d=N(\ell+1)$.
Evaluate at $X=0,\ldots,d$.
Each resulting integer matrix has entries of
polynomial bit length, and its determinant can
be computed by exact rational elimination
in polynomial time.
The determinant-size bound
$N!(\max|a_{ij}|)^N$ bounds all these outputs.
Rational interpolation at the $d+1$ distinct
integers recovers the resultant in polynomial
time; the factorial denominators of the
Lagrange basis have polynomial bit length.
Finally division by the monic polynomial $H^2$
uses polynomially many exact integer operations.
At each division step the coefficient bit bound
increases by at most the coefficient bound for
$H^2$ plus a logarithm of its degree, so it
remains polynomial throughout.
The remainder vanishes exactly when (E17.4)
holds. This proves the stated complexity
and both directions of the test.
\end{proof}

The only analytic characters in this proof are the
quadratic character used in (E17.1), the three
quadratic characters of the displayed biquadratic
field, and the trivial character.
In particular $M$, ring-class Hecke RH, general
Artin holomorphy, and elliptic $\mathrm{GL}_2$ RH
are not additional inputs.
This is a new formulation and input reduction,
with no claim of literature priority.

\paragraph{NT-E18: translation to an elliptic curve over a varying number field.}
Suppose a number field
$F=\mathbf Q[t]/(f(t))$ is given by a monic
irreducible $f\in\mathbf Q[t]$, and a nonsingular
Weierstrass equation for $E/F$ is given with
its coefficients represented in the power basis.
The defining polynomial is supplied as a dense
coefficient list, and field elements as dense
power-basis vectors. All rational coefficients
are encoded in binary, so $n=\deg f$ is at most
the total input length;
the assertion that $f$ defines a field is the
input promise.
Under $A_{\rm quad}$, geometric CM of $E$
can be decided deterministically in polynomial
time in this total input length.
No maximal order, conductor, integer factorization,
or endomorphism basis is required.

\begin{proof}
Compute $\Delta\ne0$ and $j=c_4^3/\Delta$ using
exact arithmetic in $F$.
Multiplication is reduction modulo $f$;
inversion is rational linear algebra.
These operations have polynomial bit complexity.
Let $n=\deg f$. Compute the coordinate vectors
of $1,j,\ldots,j^n$ and find the first dependence.
It gives the monic minimal polynomial
$H_j\in\mathbf Q[X]$, of degree $h\le n$.

For clarity the coefficient-size assertion does
not assume a precomputed integral basis.
The multiplication matrix of $j$ has rational
entries of polynomial bit length in the input.
After a common denominator is cleared, its first
$n$ powers have polynomial bit length: multiplying
integer matrices increases the logarithmic
entry bound by at most the previous bound plus
the fixed matrix-entry bound and $\log n$.
Cramer's determinant bounds for the first
linear dependence therefore give coefficients
of $H_j$ of polynomial bit length.
Exact elimination consequently takes polynomial
time as well.

If $H_j$ has a nonintegral coefficient, then $j$
is not an algebraic integer and cannot be a
CM $j$-invariant; return non-CM.
Otherwise $H_j$ is a monic irreducible integer
polynomial. Apply NT-E17 to it.
The degree and padded logarithmic height of
$H_j$ are polynomially bounded in the original
input length, so the complete computation is
polynomial. The characterization of CM
$j$-invariants as the roots of the Hilbert class
polynomials proves correctness.
\end{proof}

This recognizes whether
$\operatorname{End}_{\overline F}(E)$ is $\mathbf Z$
or an imaginary quadratic order.
It does not output the order discriminant or
explicit endomorphisms.
Cremona--Sutherland Algorithm~2 does output
the discriminant, and its correctness and
termination are Theorem~30, p.~15.
Its near-quasilinear expected running time is
Theorem~35, p.~17, whose hypotheses include
the separate prime-distribution heuristic (H)
and the computational heuristics in its
reference [45, \S7.1], as well as GRH.
Neither $A_D$ nor $M$ by itself removes those
extra hypotheses.
Likewise Enge's floating-point class-polynomial
complexity, Corollary~1.3 of arXiv:cs/0601104v2,
pp.~2--3, explicitly retains a rounding-error
heuristic; it is not used in either proof above.
These distinctions concern the cited algorithms,
not an assertion that unconditional CM recognition
is impossible.

Geometrically (E17.4) detects the quadratic
endomorphism algebra of $H^1(E)$ through the
cyclic-isogeny correspondences on the moduli
line of elliptic curves. The conditional
existence of a bounded correspondence degree
comes from a finite biquadratic Frobenius cover
of an arithmetic open, in the finite-etale
dictionary of SGA~1, Expose~V, \S7.
The result distinguishes CM torus monodromy
from the non-CM case without proving a
rank-two automorphic GRH.
Work/run: \texttt{cm\_recognition\_dirichlet.tex},
\texttt{GRH\_NIGHT/number\_theory}.
